% Options for packages loaded elsewhere
\PassOptionsToPackage{unicode}{hyperref}
\PassOptionsToPackage{hyphens}{url}
\documentclass[
]{article}
\usepackage{xcolor}
\usepackage[margin=1in]{geometry}
\usepackage{amsmath,amssymb}
\setcounter{secnumdepth}{-\maxdimen} % remove section numbering
\usepackage{iftex}
\ifPDFTeX
  \usepackage[T1]{fontenc}
  \usepackage[utf8]{inputenc}
  \usepackage{textcomp} % provide euro and other symbols
\else % if luatex or xetex
  \usepackage{unicode-math} % this also loads fontspec
  \defaultfontfeatures{Scale=MatchLowercase}
  \defaultfontfeatures[\rmfamily]{Ligatures=TeX,Scale=1}
\fi
\usepackage{lmodern}
\ifPDFTeX\else
  % xetex/luatex font selection
\fi
% Use upquote if available, for straight quotes in verbatim environments
\IfFileExists{upquote.sty}{\usepackage{upquote}}{}
\IfFileExists{microtype.sty}{% use microtype if available
  \usepackage[]{microtype}
  \UseMicrotypeSet[protrusion]{basicmath} % disable protrusion for tt fonts
}{}
\makeatletter
\@ifundefined{KOMAClassName}{% if non-KOMA class
  \IfFileExists{parskip.sty}{%
    \usepackage{parskip}
  }{% else
    \setlength{\parindent}{0pt}
    \setlength{\parskip}{6pt plus 2pt minus 1pt}}
}{% if KOMA class
  \KOMAoptions{parskip=half}}
\makeatother
\usepackage{color}
\usepackage{fancyvrb}
\newcommand{\VerbBar}{|}
\newcommand{\VERB}{\Verb[commandchars=\\\{\}]}
\DefineVerbatimEnvironment{Highlighting}{Verbatim}{commandchars=\\\{\}}
% Add ',fontsize=\small' for more characters per line
\usepackage{framed}
\definecolor{shadecolor}{RGB}{248,248,248}
\newenvironment{Shaded}{\begin{snugshade}}{\end{snugshade}}
\newcommand{\AlertTok}[1]{\textcolor[rgb]{0.94,0.16,0.16}{#1}}
\newcommand{\AnnotationTok}[1]{\textcolor[rgb]{0.56,0.35,0.01}{\textbf{\textit{#1}}}}
\newcommand{\AttributeTok}[1]{\textcolor[rgb]{0.13,0.29,0.53}{#1}}
\newcommand{\BaseNTok}[1]{\textcolor[rgb]{0.00,0.00,0.81}{#1}}
\newcommand{\BuiltInTok}[1]{#1}
\newcommand{\CharTok}[1]{\textcolor[rgb]{0.31,0.60,0.02}{#1}}
\newcommand{\CommentTok}[1]{\textcolor[rgb]{0.56,0.35,0.01}{\textit{#1}}}
\newcommand{\CommentVarTok}[1]{\textcolor[rgb]{0.56,0.35,0.01}{\textbf{\textit{#1}}}}
\newcommand{\ConstantTok}[1]{\textcolor[rgb]{0.56,0.35,0.01}{#1}}
\newcommand{\ControlFlowTok}[1]{\textcolor[rgb]{0.13,0.29,0.53}{\textbf{#1}}}
\newcommand{\DataTypeTok}[1]{\textcolor[rgb]{0.13,0.29,0.53}{#1}}
\newcommand{\DecValTok}[1]{\textcolor[rgb]{0.00,0.00,0.81}{#1}}
\newcommand{\DocumentationTok}[1]{\textcolor[rgb]{0.56,0.35,0.01}{\textbf{\textit{#1}}}}
\newcommand{\ErrorTok}[1]{\textcolor[rgb]{0.64,0.00,0.00}{\textbf{#1}}}
\newcommand{\ExtensionTok}[1]{#1}
\newcommand{\FloatTok}[1]{\textcolor[rgb]{0.00,0.00,0.81}{#1}}
\newcommand{\FunctionTok}[1]{\textcolor[rgb]{0.13,0.29,0.53}{\textbf{#1}}}
\newcommand{\ImportTok}[1]{#1}
\newcommand{\InformationTok}[1]{\textcolor[rgb]{0.56,0.35,0.01}{\textbf{\textit{#1}}}}
\newcommand{\KeywordTok}[1]{\textcolor[rgb]{0.13,0.29,0.53}{\textbf{#1}}}
\newcommand{\NormalTok}[1]{#1}
\newcommand{\OperatorTok}[1]{\textcolor[rgb]{0.81,0.36,0.00}{\textbf{#1}}}
\newcommand{\OtherTok}[1]{\textcolor[rgb]{0.56,0.35,0.01}{#1}}
\newcommand{\PreprocessorTok}[1]{\textcolor[rgb]{0.56,0.35,0.01}{\textit{#1}}}
\newcommand{\RegionMarkerTok}[1]{#1}
\newcommand{\SpecialCharTok}[1]{\textcolor[rgb]{0.81,0.36,0.00}{\textbf{#1}}}
\newcommand{\SpecialStringTok}[1]{\textcolor[rgb]{0.31,0.60,0.02}{#1}}
\newcommand{\StringTok}[1]{\textcolor[rgb]{0.31,0.60,0.02}{#1}}
\newcommand{\VariableTok}[1]{\textcolor[rgb]{0.00,0.00,0.00}{#1}}
\newcommand{\VerbatimStringTok}[1]{\textcolor[rgb]{0.31,0.60,0.02}{#1}}
\newcommand{\WarningTok}[1]{\textcolor[rgb]{0.56,0.35,0.01}{\textbf{\textit{#1}}}}
\usepackage{graphicx}
\makeatletter
\newsavebox\pandoc@box
\newcommand*\pandocbounded[1]{% scales image to fit in text height/width
  \sbox\pandoc@box{#1}%
  \Gscale@div\@tempa{\textheight}{\dimexpr\ht\pandoc@box+\dp\pandoc@box\relax}%
  \Gscale@div\@tempb{\linewidth}{\wd\pandoc@box}%
  \ifdim\@tempb\p@<\@tempa\p@\let\@tempa\@tempb\fi% select the smaller of both
  \ifdim\@tempa\p@<\p@\scalebox{\@tempa}{\usebox\pandoc@box}%
  \else\usebox{\pandoc@box}%
  \fi%
}
% Set default figure placement to htbp
\def\fps@figure{htbp}
\makeatother
\setlength{\emergencystretch}{3em} % prevent overfull lines
\providecommand{\tightlist}{%
  \setlength{\itemsep}{0pt}\setlength{\parskip}{0pt}}
\usepackage{bookmark}
\IfFileExists{xurl.sty}{\usepackage{xurl}}{} % add URL line breaks if available
\urlstyle{same}
\hypersetup{
  pdftitle={Power for Equivalence (TOST)},
  hidelinks,
  pdfcreator={LaTeX via pandoc}}

\title{Power for Equivalence (TOST)}
\author{}
\date{\vspace{-2.5em}2026-04-17}

\begin{document}
\maketitle

\subsection{Introduction}\label{introduction}

Power analysis for equivalence testing via the two-one-sided-tests
(TOST) procedure requires that both one-sided null hypotheses --- that
the true effect lies below the lower equivalence bound and that it lies
above the upper equivalence bound --- be rejected at level \(\alpha\).
The required sample size depends on the equivalence margin (set by
clinical or regulatory convention), the anticipated true effect
(typically zero or near-zero for an honest equivalence claim), and the
variance. TOST sample-size calculations are widely used in
bioequivalence studies of generic drugs (where the regulatory margin is
fixed at 80--125 \% on the log scale), in non-inferiority trials
reframed as equivalence, and in therapeutic-equivalence studies of
competing devices or formulations.

\subsection{Prerequisites}\label{prerequisites}

A working understanding of TOST equivalence testing, the role of the
equivalence margin, and the difference between equivalence (CI within
the margin) and traditional superiority testing (CI excluding the null).

\subsection{Theory}\label{theory}

For two groups with true mean difference \(\delta\) and pooled SD
\(\sigma\), and symmetric equivalence margin \(\pm \Delta\), the power
to conclude equivalence is approximately

\[\Pr\!\left(\frac{|\bar X_1 - \bar X_2 - \delta|}{\sigma \sqrt{2/n}} < z_{1-\alpha} \cdot \frac{\Delta - |\delta|}{\Delta \cdot \sigma \sqrt{2/n}}\right).\]

Two important features fall out of this expression. First, when the true
effect is at the centre of the equivalence margin (\(\delta = 0\)), the
required sample size is minimised. Second, when the true effect equals
the boundary (\(|\delta| = \Delta\)), equivalence cannot be established
regardless of \(n\) --- the test is fundamentally underpowered at the
boundary.

\subsection{Assumptions}\label{assumptions}

The equivalence margin is pre-specified and substantively justified, the
anticipated true effect \(\delta\) is realistically estimated, the data
are approximately Normal or the sample is large enough for the central
limit theorem, and the variance is known or accurately estimated from
pilot data.

\subsection{R Implementation}\label{r-implementation}

\begin{Shaded}
\begin{Highlighting}[]
\FunctionTok{library}\NormalTok{(TOSTER); }\FunctionTok{library}\NormalTok{(PowerTOST)}

\FunctionTok{sampleN.TOST}\NormalTok{(}\AttributeTok{alpha =} \FloatTok{0.05}\NormalTok{, }\AttributeTok{targetpower =} \FloatTok{0.80}\NormalTok{,}
             \AttributeTok{theta0 =} \FloatTok{1.00}\NormalTok{, }\AttributeTok{theta1 =} \FloatTok{0.80}\NormalTok{, }\AttributeTok{theta2 =} \FloatTok{1.25}\NormalTok{,}
             \AttributeTok{CV =} \FloatTok{0.20}\NormalTok{, }\AttributeTok{design =} \StringTok{"2x2"}\NormalTok{)}

\FunctionTok{power\_t\_TOST}\NormalTok{(}\AttributeTok{alpha =} \FloatTok{0.05}\NormalTok{, }\AttributeTok{low\_eqbound =} \SpecialCharTok{{-}}\DecValTok{2}\NormalTok{, }\AttributeTok{high\_eqbound =} \DecValTok{2}\NormalTok{,}
             \AttributeTok{mu =} \DecValTok{0}\NormalTok{, }\AttributeTok{sd =} \DecValTok{5}\NormalTok{, }\AttributeTok{type =} \StringTok{"two.sample"}\NormalTok{, }\AttributeTok{n =} \DecValTok{65}\NormalTok{)}
\end{Highlighting}
\end{Shaded}

\subsection{Output \& Results}\label{output-results}

\texttt{sampleN.TOST()} returns the required sample size for
bioequivalence with the regulatory 80--125 \% margin on the log scale;
for 20 \% within-subject CV and true ratio 1.00, \(n \approx 24\) in a 2
× 2 crossover. \texttt{power\_t\_TOST()} computes power for a generic
two-sample equivalence test with arbitrary symmetric or asymmetric
margin: at \(\pm 2\) unit margin, true difference 0, and SD 5,
\(n = 65\) per group achieves 81 \% power.

\subsection{Interpretation}\label{interpretation}

A reporting sentence: ``To establish bioequivalence (regulatory 80--125
\% equivalence margin on the log scale) with 80 \% power at
\(\alpha = 0.05\), assuming a 20 \% within-subject coefficient of
variation and a true geometric mean ratio of 1.00, \(n = 24\) subjects
are required in a 2 × 2 crossover design. The protocol enrols 28 to
allow for 15 \% dropout. Sensitivity analyses across CV
\(\in [15 \%, 30 \%]\) yield required \(n\) from 16 to 38, and across
true ratio \(\in [0.95, 1.05]\) from 24 to 33.'' Always state the margin
and the assumed true effect.

\subsection{Practical Tips}\label{practical-tips}

\begin{itemize}
\tightlist
\item
  The equivalence margin is a substantive and regulatory choice, not a
  statistical one; set it based on clinical or regulatory conventions
  (e.g., the 80--125 \% bioequivalence margin) or on a pre-specified
  clinical-acceptance criterion, and justify the choice in the protocol.
\item
  Power drops sharply as the anticipated true effect approaches the
  equivalence margin; always include sensitivity analysis over a
  plausible range of true effects to demonstrate that the design is
  robust to mild deviations from the central assumption.
\item
  Crossover bioequivalence trials use within-subject variance
  (\(\mathrm{CV}_w\)), which is typically smaller than between-subject
  variance; plan accordingly with the appropriate within-subject CV from
  pilot or literature.
\item
  Asymmetric equivalence margins are supported by most TOST tools
  (\texttt{low\_eqbound} ≠ \texttt{high\_eqbound}) and are useful when
  only one direction of effect carries a clinical penalty; for example,
  a drug that is allowed to be more effective but not less effective
  than the reference.
\item
  Non-inferiority testing is mathematically a one-sided TOST against a
  single non-inferiority margin; the required sample size is roughly
  half of a two-sided equivalence test for the same margin and
  assumptions, reflecting the simpler hypothesis.
\item
  For highly variable drugs (within-subject CV above 30 \%),
  reference-scaled bioequivalence widens the margin proportionally to
  the reference variability; standard fixed-margin TOST sample-size
  tools under-estimate the required \(n\) for these compounds.
\end{itemize}

\subsection{R Packages Used}\label{r-packages-used}

\texttt{PowerTOST} for canonical bioequivalence and equivalence
sample-size and power across crossover and parallel-group designs;
\texttt{TOSTER} for general TOST analysis and power on raw and
standardised effect sizes; \texttt{MBESS} and
\texttt{pwrss::pwrss.t.2means()} for related precision and equivalence
calculations; \texttt{bear} for end-to-end FDA-compliant bioequivalence
workflow; \texttt{replicateBE} for replicate-design reference-scaled
bioequivalence.

\subsection{For Reviewers}\label{for-reviewers}

\textbf{What to look for in a paper using this method.}

\begin{itemize}
\tightlist
\item
  \textbf{Common misapplications.}
\item
  \textbf{Diagnostics that should be reported but often aren't.}
\item
  \textbf{Red flags in tables and figures.}
\item
  \textbf{What to verify.}
\item
  \textbf{What an adequate Methods paragraph must contain.}
\end{itemize}

\subsection{See also --- labs in this
chapter}\label{see-also-labs-in-this-chapter}

\begin{itemize}
\tightlist
\item
  \href{labs/lab-sample-size-power-and-quarto-reporting.Rmd}{Sample
  size, power, and Quarto reporting}
\item
  \href{labs/lab-power-closed-form-and-simulation.Rmd}{Power:
  closed-form and simulation}
\end{itemize}

\end{document}
